\section{Conclusion}
We have shown that generating Wikipedia can be approached as a multi-document
summarization problem with a large, parallel dataset,
and demonstrated a two-stage extractive-abstractive framework for carrying it out.
The coarse extraction method used in the first
stage appears to have a significant effect on final performance, suggesting
further research on improving it would be fruitful. We introduce a new,
decoder-only sequence transduction model for the abstractive stage,
capable of handling very long input-output examples.
This model significantly outperforms traditional encoder-decoder architectures
on long sequences, allowing us to condition on many reference documents and
to generate coherent and informative Wikipedia articles.
